\section{Universal properties of the two rigidifications}
\label{sec:rigidification-v151}

The effective stacks used above have a universal characterization.
This specifies what is forgotten from isomorphisms, as opposed to what
is retained from objects.  We use rigidification in its standard sense
of removing a flat normal subgroup of inertia; see
\cite[Tag~04V2]{StacksRigidV151}.  The argument below gives the
universal property directly in the present situation, including the
nonconstant subgroup obtained by descent.

Let \(H_0\subset\Gr(r,S_D(E))\) be a
\(\operatorname{GL}(E)\)-invariant locally closed subscheme.  Let
\(\mathscr A_{H_0}\) be the stack of algebras satisfying the trace,
filtration and Hilbert-function conditions of
Proposition~\ref{prop:relative-first-relation}, whose intrinsic
kernel is in \(H_0\).  Write
\(\mathscr R_{H_0}=[H_0/\operatorname{GL}(E)]\).

\begin{theorem}[First-relation rigidification]
\label{thm:rigidification-v151}
Taking the first relation defines an fppf gerbe
\[
                      q:\mathscr A_{H_0}\longrightarrow\mathscr R_{H_0}.
\]
Its relative inertia \(\mathscr U\) consists exactly of the
substitutions which induce the identity on \(\mathcal N/\mathcal N^2\).
It is a smooth affine normal subgroup of the inertia.  In a homogeneous
presentation its underlying scheme is
\(\Hom(\mathcal E,\mathcal N^2)\), with the substitution group law.
For every stack \(\mathscr Z\) in groupoids on complex schemes,
precomposition with \(q\) identifies
\begin{equation}\label{eq:rigidification-universal-v151}
 \Hom(\mathscr R_{H_0},\mathscr Z)
 \simeq
 \{F\in\Hom(\mathscr A_{H_0},\mathscr Z):F(\mathscr U)=1\}.
\end{equation}
The right side is the full subgroupoid, including all natural
isomorphisms, of functors killing the indicated relative inertia.
The construction is compatible with base change on the relation stack.
\end{theorem}
\begin{proof}
A relation bundle has its homogeneous truncated algebra as a lift.
Every algebra over that relation bundle is locally isomorphic to this
lift.  The first-relation proposition shows that an isomorphism of
relation bundles lifts fppf locally, and any two lifts differ by a
unique element of the indicated kernel.  These assertions are exactly
the gerbe and relative-inertia assertions.  The kernel is intrinsic,
so it is stable under conjugation.  Its local affine-space description
shows smoothness and flatness; no additive group law is being asserted.

For clarity, we prove the descent implicit in
\eqref{eq:rigidification-universal-v151}.  Given \(F\) on the right
and an object \(y\) of \(\mathscr R_{H_0}\), choose local lifts
\(a_i\) over an fppf cover.  Locally choose comparison isomorphisms
on overlaps above the identity of \(y\).  Two choices differ by
\(\mathscr U\), hence have the same image under \(F\).  On triple
overlaps the failure of the chosen lifts to satisfy the cocycle
condition is again in \(\mathscr U\).  Their images therefore give
a genuine descent datum for \(F(a_i)\), which is effective because
\(\mathscr Z\) is a stack.  Refinements and different lift choices
give canonically isomorphic descended objects.  An arrow between two
objects of \(\mathscr R_{H_0}\) is treated in the same way, using
local lifts of that arrow.  This constructs the descended functor.
Natural isomorphisms descend by the same argument, proving full
faithfulness as well as essential surjectivity.  Conversely every
functor pulled back by \(q\) kills its relative inertia.  All the
constructions use covers and isomorphism sheaves, so commute with
base change.
\end{proof}

\begin{corollary}[The effective pencil stack is a successive rigidification]
\label{cor:two-rigidifications-v151}
On the pencil relation stratum the second morphism
\[
 [\mathscr P/G^\circ]\longrightarrow[\mathscr P/\PGL(V)],
                     \qquad \mathscr P=\Gr(2,\Sym^2V),
\]
is the rigidification by the normal relative inertia induced from
\(\ker(G^\circ\to\PGL(V))=\operatorname{GL}(U)\).
Consequently Theorem~\ref{thm:pencil-stack-equivalence} has the
universal property of first killing \(\mathscr U\) and then this
right-factor inertia.  It does not kill the stabilizer of the pencil
in \(\PGL(V)\).
\end{corollary}
\begin{proof}
The right factor acts trivially on \(\mathscr P\), and the quotient
map of groups is fppf surjective.  Objects and isomorphisms of the
second quotient therefore lift locally to the first, with differences
in that kernel.  Apply the descent proof of
Theorem~\ref{thm:rigidification-v151}.  Over an object represented by
a \(G^\circ\)-torsor the subgroup is its conjugation-associated
\(\operatorname{GL}(U)\)-bundle; it need not be a constant group
scheme on the base.  This describes the actual subgroup stack and
shows that no global splitting is assumed.  Composing with the first
rigidification gives the asserted successive universal property.
All remaining pencil automorphisms survive in the target.
\end{proof}
